\subsection{Topology, contracts and the Native result}\label{sec:native-gap}
The reading guide below applies only to the retained full-matrix figures, not the compact run-9 topology.
\textbf{How to read these figures.} Native: one process, one SSD. Tidehunter: five owners. Other external contributors record five configured or observed storage nodes; Cassandra's 3 November run records five contacts, not a verified five-node ring. Five storage-host SSD mounts are attested only for the inspected 4 November and final deployments. FoundationDB also has a separate log/coordinator host. Equal RF1 is not equal SSD count. Native performs file and directory fsync after each put; owner covering fdatasync can serve multiple batches. Valkey waits for AOF fsync once per event-loop pass\footnote{Valkey 9.0.7, \texttt{src/server.c}: \texttt{beforeSleep} flushes AOF before pending client writes. See \url{https://github.com/valkey-io/valkey/blob/9.0.7/src/server.c} and \url{https://valkey.io/topics/persistence/}. Commands per sync were not measured in these runs.}; Kvrocks uses a synced RocksDB WAL; Garnet waits for AOF commit. Hollow write markers mean Cassandra, FoundationDB and TiKV device persistence is not verified. \runs. Hardware and configuration are not fully matched.

API time is completion minus dispatch; response time is completion minus schedule time, before dispatch and admission. Both use successful calls. Valkey reads RAM; Native uses a resident-oriented read workload, with per-call source not established; owner and hybrid cache paths differ. In the two c16 read figures, panel four is mean in-flight successful API calls: each run's sum of successful API times divided by elapsed case time, then the median across external runs. In the two c1 figures, it is the selected operation's API mean. Low occupancy means that finite rates do not establish store capacity. Both axes are logarithmic; x-axis K means binary bytes, while y-axis k and M are decimal.
\textbf{Topology and durable mechanics.} Native is one process on one SSD; the retained final cases of other systems use five storage nodes. Native's retained per-put benchmark barrier calls \texttt{Db::insert}, then fsyncs every registered Value-WAL file and opens and fsyncs the directory after each put. That arm does not group scalar callers; the separate run-9 grouped Native arm does. Owner synchronous acknowledges after its covering file fdatasync; batches appended during one sync can share the next sync. A missing network hop does not make Native a latency lower bound under another contract or topology. These settings express durable intent, not proved crash or device-persistence equivalence.

\begin{table}[H]\centering\fitinput{tab-per-node.tex}
\caption{64 KiB puts, c16: aggregate rate / retained-final node count. Even spread assumed for callers, rate and payload; these are normalizations, not observed node-level metrics. Native: one process, one SSD. Tidehunter: five owners. Other external contributors record five configured or observed storage nodes; Cassandra's 3 November run records five contacts, not a verified five-node ring. Five storage-host SSD mounts are attested only for the inspected 4 November and final deployments. FoundationDB excludes its separate log/coordinator host. Rates use finite case time, including work outside API calls. Native: per-put file/directory barrier; TH: covering sync; Valkey: AOF-always; Garnet: AOF commit; Kvrocks: synced WAL. Other writes have unverified device persistence. No cache drop; coordinator cache off. \runs. c1 is mixed, not pure-write.}\label{tab:pernode}
\end{table}
Final counts come from case configs and observed topology, recorded in \path{storage-topology.csv}. Other external contributors record five configured or observed storage nodes; Cassandra's 3 November run records five contacts, not a verified five-node ring (\path{historical-topology.csv}). This is not verification from every case configuration or proof of five physical SSDs. Five storage-host SSD mounts are attested only for the inspected 4 November and final deployments. FoundationDB excludes its separate log/coordinator host; TiKV excludes PD. Some 3 November FoundationDB machine ids and TiKV host labels name a different fleet from the report path; the path does not establish execution-fleet identity.

The gate-window completed-write counts span \N{B:device-write-spread:th_verbs:minimum|,.0f}--\N{B:device-write-spread:th_verbs:maximum|,.0f} across owner verbs nodes (max/min \N{B:device-write-spread:th_verbs:ratio|.2f}) and \N{B:device-write-spread:valkey_rdma:minimum|,.0f}--\N{B:device-write-spread:valkey_rdma:maximum|,.0f} across Valkey RDMA primaries (\N{B:device-write-spread:valkey_rdma:ratio|.2f}). These final-run counters include background work and use different node windows; they are not API-linked rates or the historical median case.

\textbf{The Native observation.} At 64 KiB, barrier c16 p50/mean are \N{B:per-node-64k:native:p50_us}/\N{B:per-node-64k:native:mean_us}~\us, with \N{B:per-node-64k:native:puts_s|,.0f} puts/s and \N{B:per-node-64k:native:mean_inflight_calls} in-flight calls. Mixed c1 barrier p50/mean are \N{p50:native:primary_rf1:point-mixed:put:65536|.1f}/\N{mean:native:primary_rf1:point-mixed:put:65536|.1f}~\us. The mean ratio is \N{J:native_mean_ratio}; the reciprocal c1 API mean is \N{J:native_reciprocal|,.0f} puts/s, not measured c1 throughput, and c16 reaches \N{J:native_reciprocal_pct|.0f} percent of that reference. \inference\ The illustration assumes that mixed-stream c1 API service represents pure-write service; the workloads differ, so the ratios do not demonstrate serialization. A serialized-service model with one independent barrier per put is consistent with the ratios, but API calls overlap; no trace separates sync serialization, queue depth, faults, allocation or journal work.

Native page-cache acknowledgement has no per-put foreground fsync or per-put persistence guarantee. Its 64 KiB c16 p50 is \N{J:pagecache64}~\us\ and mixed c1 p50 is \N{J:pagecache_c1_64}~\us. Its c16 p50 is \N{J:pagecache_jump} times the 8 KiB value of \N{J:pagecache8}~\us. The cause is not established. \inference\ Faults, allocation, writeback and fragment events are candidates, not measured contributions. The rate panel omits this page-cache series: its API occupancy is low through 8 KiB but higher at 64 KiB. Owner page-cache write curves were not run on this fleet and are not drawn.

\textbf{Where Valkey is ahead, and where it is not.} The c1 mixed-stream readings and c16 figures show different populations, not a controlled concurrency sweep. At 1 KiB, owner synchronous verbs is ahead of historical Valkey RDMA AOF-always on c1 put p50; owner resident-oriented reads are also ahead of Valkey RAM GET at c1. At c16, the 1 KiB read p50 values are close (owner \N{p50:th_verbs:primary_rf1:read-concurrent:read:1024|.1f}~\us; Valkey \N{p50:valkey_rdma:primary_rf1:read-concurrent:read:1024|.1f}~\us). Valkey is ahead on c16 put p50 at 1 and 64 KiB, and c1 mixed put at 64 KiB; owner read p50 is ahead at 64 KiB for c1 and c16. This is not a claim that every store wins at every size or quantile: Garnet's commit-wait c16 put p99 at 1 KiB is \N{p99ms:garnet:primary_rf1:write-concurrent:put:1024|.1f}~ms; the stall cause is not established.

At 64 KiB c16, the even-spread rate is \N{B:per-node-64k:valkey_rdma:puts_s_node|,.0f} puts/s/node for Valkey RDMA AOF-always, \N{B:per-node-64k:th_verbs:puts_s_node|,.0f} for owner synchronous verbs, and \N{B:per-node-64k:native:puts_s_node|,.0f} for the one-SSD Native barrier. Group commit is possible in Valkey and the owner; actual commands per sync are not measured. The aggregate rate alone cannot identify batch size, sync time or the cause of their difference.

\inference\ Waits, hand-offs and size-dependent software are plausible owner limits, not proof that the SSD has no role. On the separate run-8 fleet, synchronous verbs c16 1 KiB has \N{W:p16v|.0f} percent of its same-case API mean in three waits (\N{C:p16v:coalescing}, \N{C:p16v:queue}, \N{C:p16v:durable} of \N{D:p16v:mean_us}~\us). At c1 64 KiB, coordinator plus owner software is \N{J:owner64_software_pct|.0f} percent (\N{C:p64v:coordinator} plus \N{C:p64v:owner} of \N{D:p64v:mean_us}~\us). These intervals are not avoidable-cost bounds and must not be subtracted from separately sampled Valkey measurements.

\textbf{Fleet and profile context.} Run-8 unprofiled c1 1 KiB synchronous verbs put p50/mean is \N{D:uput1v:p50_us}/\N{D:uput1v:mean_us}~\us; the full-matrix mixed point put is \N{p50:th_verbs:primary_rf1:point-mixed:put:1024|.1f}/\N{mean:th_verbs:primary_rf1:point-mixed:put:1024|.1f}~\us. They differ in fleet, workload and repetition count. This difference is not a measured regression or profiling cost. Run 8 also compares one Native SSD with five owners; its c16 64 KiB Native barrier is not an equal-device-count control.
The compact one-owner control in Section~\ref{sec:rwc} is a distinct run, not a before/after regression test or a replacement for these five-owner populations.

\textbf{Unexplained features.} Garnet AOF-commit c16 put p99 stays near \N{p99ms:garnet:primary_rf1:write-concurrent:put:1024|.1f}--\N{p99ms:garnet:primary_rf1:write-concurrent:put:256|.1f}~ms across sizes; retained 64 B samples show stalls of about \N{B:chart-anomalies:garnet:stall_min_ms|.0f}--\N{B:chart-anomalies:garnet:stall_max_ms|.0f}~ms. FoundationDB p50 plateaus near \N{p50:foundationdb:primary_rf1:write-concurrent:put:64|.0f}~\us, then dips to \N{p50:foundationdb:primary_rf1:write-concurrent:put:8192|.0f}~\us\ at 8 KiB; Cassandra steps from \N{p50:cassandra:primary_rf1:write-concurrent:put:256|.0f} to \N{p50:cassandra:primary_rf1:write-concurrent:put:1024|.0f}~\us\ between 256 B and 1 KiB. Both have unverified device persistence. Native barrier p50 dips from \N{p50:native:primary_rf1:write-concurrent:put:256|.1f} to \N{p50:native:primary_rf1:write-concurrent:put:1024|.1f}~\us. Owner synchronous kernel/SPDK put curves lie above TCP/verbs. No cause is claimed. See \path{external-spread.csv} for external per-run spread; no whiskers are drawn.

